\documentclass[10pt,aspectratio=169]{beamer}

\usetheme{Madrid}
\usecolortheme{seahorse}
\setbeamertemplate{navigation symbols}{}
\setbeamertemplate{caption}[numbered]

\usepackage[utf8]{inputenc}
\usepackage[T1]{fontenc}
\usepackage[spanish,es-tabla]{babel}
\usepackage{amsmath, amssymb, booktabs, array}
\usepackage{tikz}
\usepackage{xcolor}
\usepackage{listings}

\definecolor{codebg}{RGB}{248,248,248}
\definecolor{codekw}{RGB}{30,70,140}
\definecolor{codestr}{RGB}{140,70,30}
\definecolor{codecom}{RGB}{80,120,80}
\definecolor{darkblue}{RGB}{25,55,95}

\lstdefinestyle{Rstyle}{
  language=R,
  basicstyle=\ttfamily\scriptsize,
  keywordstyle=\color{codekw}\bfseries,
  commentstyle=\color{codecom},
  stringstyle=\color{codestr},
  backgroundcolor=\color{codebg},
  showstringspaces=false,
  breaklines=true,
  frame=single,
  rulecolor=\color{gray!40},
  columns=fullflexible,
  keepspaces=true
}

\title[Análisis de asociación]{Análisis de asociación}
\subtitle{Reglas de asociación, algoritmos clásicos y ejemplos en R}
\author{Minería de Datos}
\institute{Clase práctica}
\date{\today}

\begin{document}

%------------------------------------------------
\begin{frame}
  \titlepage
\end{frame}

%------------------------------------------------
\begin{frame}{Objetivos de la clase}
\begin{itemize}
  \item Comprender el problema de descubrir asociaciones frecuentes en datos transaccionales.
  \item Definir itemsets, reglas, soporte, confianza, lift y otras medidas de interés.
  \item Revisar los algoritmos clásicos: Apriori, Eclat y FP-Growth.
  \item Implementar un flujo de trabajo en R usando \texttt{arules} y \texttt{arulesViz}.
  \item Interpretar reglas de asociación y discutir sus limitaciones.
\end{itemize}
\end{frame}

%------------------------------------------------
\begin{frame}{Motivación: de datos a patrones}
\begin{columns}[T]
\column{0.52\textwidth}
\textbf{Pregunta central:}\vspace{0.2cm}

¿Qué elementos tienden a aparecer juntos más frecuentemente de lo esperado?

\vspace{0.3cm}
\begin{itemize}
  \item Canasta de compras: pan, leche, mantequilla.
  \item Diagnóstico: síntomas, exámenes y enfermedad.
  \item Web: páginas visitadas en una misma sesión.
  \item Biología: genes, mutaciones o rutas alteradas en conjunto.
\end{itemize}

\column{0.42\textwidth}
\begin{block}{Idea}
Una regla de asociación resume una relación condicional frecuente:
\[
A \Rightarrow B
\]
``si aparece A, también suele aparecer B''.
\end{block}
\end{columns}
\end{frame}

%------------------------------------------------
\begin{frame}{Datos transaccionales}
\begin{table}
\centering
\begin{tabular}{cl}
\toprule
Transacción & Ítems observados \\
\midrule
T1 & pan, leche, huevos \\
T2 & pan, mantequilla \\
T3 & leche, huevos, cereal \\
T4 & pan, leche, mantequilla \\
T5 & pan, leche, huevos, mantequilla \\
\bottomrule
\end{tabular}
\end{table}

\vspace{0.3cm}
Cada fila representa una unidad de análisis; cada ítem puede estar presente o ausente. El formato usual es una matriz binaria dispersa: muchas columnas, muchos ceros.
\end{frame}

%------------------------------------------------
\begin{frame}{Conceptos básicos}
\begin{block}{Ítem e itemset}
Un \textbf{ítem} es una característica discreta. Un \textbf{itemset} es un conjunto de ítems, por ejemplo \{pan, leche\}.
\end{block}

\begin{block}{Regla de asociación}
Una regla se escribe como:
\[
X \Rightarrow Y, \quad X \cap Y = \emptyset
\]
Donde $X$ es el antecedente y $Y$ el consecuente.
\end{block}

\begin{exampleblock}{Ejemplo}
\[
\{pan, mantequilla\} \Rightarrow \{leche\}
\]
Interpretación: entre las transacciones con pan y mantequilla, una fracción relevante también contiene leche.
\end{exampleblock}
\end{frame}

%------------------------------------------------
\begin{frame}{Soporte}
\textbf{Soporte de un itemset:} proporción de transacciones que contienen ese conjunto.
\[
\text{support}(X) = \frac{\#(X)}{N}
\]

\vspace{0.3cm}
\textbf{Soporte de una regla:}
\[
\text{support}(X \Rightarrow Y) = \text{support}(X \cup Y)
\]

\begin{block}{Interpretación}
El soporte mide frecuencia absoluta relativa. Es útil para filtrar patrones demasiado raros, aunque patrones raros pueden ser importantes en algunos dominios.
\end{block}
\end{frame}

%------------------------------------------------
\begin{frame}{Confianza}
\textbf{Confianza:} probabilidad empírica de observar $Y$ dado que se observa $X$.
\[
\text{confidence}(X \Rightarrow Y) = \frac{\text{support}(X \cup Y)}{\text{support}(X)}
\]

\begin{exampleblock}{Ejemplo}
Si 40\% de las transacciones tienen pan y 30\% tienen pan y leche:
\[
\text{confidence}(\{pan\} \Rightarrow \{leche\}) = \frac{0.30}{0.40} = 0.75
\]
\end{exampleblock}

\begin{alertblock}{Cuidado}
Una confianza alta no implica que la asociación sea interesante: el consecuente puede ser muy frecuente por sí solo.
\end{alertblock}
\end{frame}

%------------------------------------------------
\begin{frame}{Lift}
\textbf{Lift:} compara la co-ocurrencia observada con la esperada bajo independencia.
\[
\text{lift}(X \Rightarrow Y)=\frac{\text{confidence}(X \Rightarrow Y)}{\text{support}(Y)}
\]

\begin{itemize}
  \item $\text{lift} = 1$: X e Y aparecen juntos como se esperaría por azar.
  \item $\text{lift} > 1$: asociación positiva.
  \item $\text{lift} < 1$: asociación negativa.
\end{itemize}

\begin{block}{Idea clave}
El lift corrige parcialmente el problema de consecuentes muy frecuentes.
\end{block}
\end{frame}

%------------------------------------------------
\begin{frame}{Otras medidas de interés}
\begin{table}
\centering
\begin{tabular}{p{0.22\textwidth}p{0.62\textwidth}}
\toprule
Medida & Uso principal \\
\midrule
Leverage & Diferencia entre co-ocurrencia observada y esperada. \\
Conviction & Mide cuánto se reduce el error al usar la regla. \\
Coverage & Proporción de transacciones cubiertas por el antecedente. \\
Chi-cuadrado & Evalúa independencia entre antecedente y consecuente. \\
Fisher exacto & Alternativa para tablas pequeñas o eventos raros. \\
\bottomrule
\end{tabular}
\end{table}

\vspace{0.2cm}
En la práctica se combinan varias medidas: soporte mínimo, confianza mínima, lift alto y una interpretación coherente con el dominio.
\end{frame}

%------------------------------------------------
\begin{frame}{Propiedad antimonótona del soporte}
\begin{block}{Principio clave}
Si un itemset no es frecuente, ningún superconjunto suyo puede ser frecuente.
\end{block}

\[
\text{support}(X \cup Z) \leq \text{support}(X)
\]

\vspace{0.3cm}
\begin{itemize}
  \item Esta propiedad permite podar el espacio de búsqueda.
  \item Es la base del algoritmo Apriori.
  \item Reduce drásticamente el número de combinaciones evaluadas.
\end{itemize}
\end{frame}

%------------------------------------------------
\begin{frame}{Algoritmo Apriori: intuición}
\begin{enumerate}
  \item Encontrar todos los ítems frecuentes de tamaño 1.
  \item Generar candidatos de tamaño 2 combinando ítems frecuentes.
  \item Eliminar candidatos con subconjuntos no frecuentes.
  \item Contar soporte en la base de transacciones.
  \item Repetir para itemsets de tamaño $k$.
  \item Generar reglas desde los itemsets frecuentes.
\end{enumerate}

\begin{alertblock}{Costo computacional}
Apriori puede requerir muchas pasadas sobre los datos y generar muchos candidatos, especialmente con soporte bajo.
\end{alertblock}
\end{frame}

%------------------------------------------------
\begin{frame}{Apriori: pseudocódigo}
\begin{block}{Entrada}
Base transaccional $D$, soporte mínimo $s_{min}$, confianza mínima $c_{min}$.
\end{block}

\begin{enumerate}
  \item $L_1 =$ itemsets frecuentes de tamaño 1.
  \item Para $k=2,3,\ldots$ mientras $L_{k-1}$ no esté vacío:
  \begin{enumerate}
    \item Generar candidatos $C_k$ desde $L_{k-1}$.
    \item Podar candidatos con subconjuntos no frecuentes.
    \item Contar soporte de cada candidato en $D$.
    \item Guardar en $L_k$ los candidatos con soporte $\geq s_{min}$.
  \end{enumerate}
  \item A partir de $\cup_k L_k$, generar reglas que cumplen $c_{min}$.
\end{enumerate}
\end{frame}

%------------------------------------------------
\begin{frame}{Algoritmo Eclat}
\begin{columns}[T]
\column{0.48\textwidth}
\textbf{Idea:} representar cada ítem por la lista de transacciones en que aparece.

\vspace{0.2cm}
\begin{itemize}
  \item Usa formato vertical.
  \item El soporte de un itemset se calcula como intersección de listas.
  \item Suele ser eficiente cuando las listas son manejables.
\end{itemize}

\column{0.44\textwidth}
\begin{exampleblock}{Ejemplo}
\begin{tabular}{ll}
pan & T1, T2, T4, T5 \\
leche & T1, T3, T4, T5 \\
\end{tabular}

\vspace{0.2cm}
\[
\{pan, leche\}: T1,T4,T5
\]
Soporte $=3/5$.
\end{exampleblock}
\end{columns}
\end{frame}

%------------------------------------------------
\begin{frame}{Algoritmo FP-Growth}
\begin{itemize}
  \item Evita generar candidatos explícitos como Apriori.
  \item Comprime la base de transacciones en un árbol de patrones frecuentes: \textbf{FP-tree}.
  \item Explora bases condicionales para extraer itemsets frecuentes.
\end{itemize}

\begin{block}{Ventaja}
Puede ser mucho más rápido que Apriori en bases grandes o densas, porque reduce pasadas por los datos y evita una explosión de candidatos.
\end{block}

\begin{alertblock}{Consideración}
La implementación y la estructura de memoria son más complejas que en Apriori.
\end{alertblock}
\end{frame}

%------------------------------------------------
\begin{frame}{Comparación conceptual}
\begin{table}
\centering
\begin{tabular}{p{0.2\textwidth}p{0.24\textwidth}p{0.22\textwidth}p{0.2\textwidth}}
\toprule
Algoritmo & Representación & Estrategia & Fortalezas \\
\midrule
Apriori & Horizontal & Genera y poda candidatos & Simple, interpretable, clásico \\
Eclat & Vertical & Intersección de tidsets & Eficiente para ciertos datos dispersos \\
FP-Growth & Árbol comprimido & Crecimiento de patrones & Rápido en bases grandes/densas \\
\bottomrule
\end{tabular}
\end{table}
\end{frame}


%------------------------------------------------
\begin{frame}{Dataset común para cálculo manual}
\small
Usaremos el mismo ejemplo para ilustrar los tres algoritmos clásicos.

\vspace{0.15cm}
\begin{columns}[T]
\column{0.48\textwidth}
\begin{table}
\centering
\scriptsize
\begin{tabular}{cl}
\toprule
ID & Ítems \\
\midrule
T1 & pan, leche, huevos \\
T2 & pan, mantequilla \\
T3 & leche, huevos, cereal \\
T4 & pan, leche, mantequilla \\
T5 & pan, leche, huevos, mantequilla \\
\bottomrule
\end{tabular}
\end{table}

\column{0.45\textwidth}
\begin{block}{Umbral de trabajo}
\[
N=5, \qquad s_{min}=0.40
\]
\[
\text{soporte mínimo en conteo}=0.40 \times 5=2
\]
Un itemset es frecuente si aparece en al menos 2 transacciones.
\end{block}
\end{columns}
\end{frame}


%================================================
\section{Cálculo manual de los algoritmos clásicos}

%------------------------------------------------
\begin{frame}{Cálculo manual: dataset y umbral}
\scriptsize
\begin{columns}[T]
\column{0.48\textwidth}
\textbf{Base transaccional $D$}
\begin{table}
\centering
\begin{tabular}{cl}
\toprule
ID & Compra \\
\midrule
T1 & pan, leche, huevos \\
T2 & pan, mantequilla \\
T3 & leche, huevos, cereal \\
T4 & pan, leche, mantequilla \\
T5 & pan, leche, huevos, mantequilla \\
\bottomrule
\end{tabular}
\end{table}

\column{0.46\textwidth}
\textbf{Parámetros para todos los algoritmos}
\[
N=5, \qquad s_{min}=0.40
\]
\[
\text{conteo mínimo}=s_{min}\times N=0.40\times 5=2
\]
\begin{block}{Regla de decisión}
Un itemset se considera frecuente si aparece en \textbf{2 o más} transacciones.
\end{block}
\end{columns}
\end{frame}

%------------------------------------------------
\begin{frame}{Apriori a mano I: contar $C_1$ y obtener $L_1$}
\scriptsize
\textbf{Idea del algoritmo:} primero contar ítems individuales. Luego solo se combinan los ítems frecuentes.

\vspace{0.15cm}
\begin{columns}[T]
\column{0.47\textwidth}
\begin{table}
\centering
\begin{tabular}{lcl}
\toprule
Ítem & Conteo & Soporte \\
\midrule
pan & 4 & $4/5=0.80$ \\
leche & 4 & $4/5=0.80$ \\
huevos & 3 & $3/5=0.60$ \\
mantequilla & 3 & $3/5=0.60$ \\
cereal & 1 & $1/5=0.20$ \\
\bottomrule
\end{tabular}
\end{table}

\column{0.48\textwidth}
\textbf{Aplicando soporte mínimo}
\[
L_1=\{pan, leche, huevos, mantequilla\}
\]
\[
cereal \notin L_1 \quad \text{porque} \quad 1 < 2
\]
\begin{alertblock}{Poda Apriori}
Como cereal no es frecuente, ningún itemset que contenga cereal puede ser frecuente.
\end{alertblock}
\end{columns}
\end{frame}

%------------------------------------------------
\begin{frame}{Apriori a mano II: generar $C_2$ y obtener $L_2$}
\scriptsize
\textbf{Desde $L_1$ se generan candidatos de tamaño 2.} Luego se cuentan en las 5 transacciones.

\begin{table}
\centering
\begin{tabular}{lcll}
\toprule
Candidato $C_2$ & Conteo & Soporte & Decisión \\
\midrule
\{pan, leche\} & 3 & $3/5=0.60$ & frecuente \\
\{pan, huevos\} & 2 & $2/5=0.40$ & frecuente \\
\{pan, mantequilla\} & 3 & $3/5=0.60$ & frecuente \\
\{leche, huevos\} & 3 & $3/5=0.60$ & frecuente \\
\{leche, mantequilla\} & 2 & $2/5=0.40$ & frecuente \\
\{huevos, mantequilla\} & 1 & $1/5=0.20$ & se poda \\
\bottomrule
\end{tabular}
\end{table}

\[
L_2=\{\{pan,leche\},\{pan,huevos\},\{pan,mantequilla\},\{leche,huevos\},\{leche,mantequilla\}\}
\]
\end{frame}

%------------------------------------------------
\begin{frame}{Apriori a mano III: generar $C_3$, podar y calcular reglas}
\scriptsize
\begin{columns}[T]
\column{0.50\textwidth}
\textbf{Candidatos de tamaño 3}
\begin{table}
\centering
\begin{tabular}{lcc}
\toprule
Candidato $C_3$ & Conteo & Decisión \\
\midrule
\{pan, leche, huevos\} & 2 & queda \\
\{pan, leche, mantequilla\} & 2 & queda \\
\{pan, huevos, mantequilla\} & -- & se poda \\
\{leche, huevos, mantequilla\} & -- & se poda \\
\bottomrule
\end{tabular}
\end{table}

Los dos últimos se podan porque contienen:
\[
\{huevos,mantequilla\}\notin L_2
\]

\column{0.46\textwidth}
\textbf{Ejemplo de regla desde $L_3$}
\[
\{pan,huevos\}\Rightarrow\{leche\}
\]
\[
\text{support}=\frac{\#(pan,huevos,leche)}{5}=\frac{2}{5}=0.40
\]
\[
\text{confidence}=\frac{\#(pan,huevos,leche)}{\#(pan,huevos)}=\frac{2}{2}=1.00
\]
\[
\text{lift}=\frac{1.00}{support(leche)}=\frac{1.00}{4/5}=1.25
\]
\end{columns}
\end{frame}

%------------------------------------------------
\begin{frame}{Eclat a mano I: convertir a formato vertical}
\scriptsize
\textbf{Idea del algoritmo:} cada ítem guarda el conjunto de transacciones en que aparece.

\begin{table}
\centering
\begin{tabular}{llcc}
\toprule
Ítem & TIDset & Conteo & Decisión \\
\midrule
pan & \{T1,T2,T4,T5\} & 4 & frecuente \\
leche & \{T1,T3,T4,T5\} & 4 & frecuente \\
huevos & \{T1,T3,T5\} & 3 & frecuente \\
mantequilla & \{T2,T4,T5\} & 3 & frecuente \\
cereal & \{T3\} & 1 & se elimina \\
\bottomrule
\end{tabular}
\end{table}

\[
L_1=\{pan, leche, huevos, mantequilla\}
\]
\begin{block}{Diferencia clave con Apriori}
Eclat calcula soportes mediante \textbf{intersecciones de TIDsets}, no generando y recontando todos los candidatos sobre la tabla original.
\end{block}
\end{frame}

%------------------------------------------------
\begin{frame}{Eclat a mano II: intersecciones para obtener itemsets frecuentes}
\scriptsize
\begin{columns}[T]
\column{0.50\textwidth}
\textbf{Pares frecuentes}
\begin{align*}
T(pan)\cap T(leche)&=\{T1,T4,T5\}\Rightarrow 3 \\
T(pan)\cap T(huevos)&=\{T1,T5\}\Rightarrow 2 \\
T(pan)\cap T(mantequilla)&=\{T2,T4,T5\}\Rightarrow 3 \\
T(leche)\cap T(huevos)&=\{T1,T3,T5\}\Rightarrow 3 \\
T(leche)\cap T(mantequilla)&=\{T4,T5\}\Rightarrow 2 \\
T(huevos)\cap T(mantequilla)&=\{T5\}\Rightarrow 1
\end{align*}
\[
\{huevos,mantequilla\}\text{ se elimina}
\]

\column{0.45\textwidth}
\textbf{Triples por intersección}
\begin{align*}
T(pan,leche)\cap T(huevos)
&=\{T1,T4,T5\}\cap\{T1,T3,T5\}\\
&=\{T1,T5\}\Rightarrow 2
\end{align*}
\[
\{pan,leche,huevos\}\text{ es frecuente}
\]
\begin{align*}
T(pan,leche)\cap T(mantequilla)
&=\{T1,T4,T5\}\cap\{T2,T4,T5\}\\
&=\{T4,T5\}\Rightarrow 2
\end{align*}
\[
\{pan,leche,mantequilla\}\text{ es frecuente}
\]
\end{columns}
\end{frame}

%------------------------------------------------
\begin{frame}[fragile]{FP-Growth a mano I: ordenar transacciones frecuentes}
\scriptsize
\textbf{Idea del algoritmo:} comprimir la base en un FP-tree. Primero se eliminan ítems no frecuentes y se ordenan por frecuencia global.

\begin{columns}[T]
\column{0.41\textwidth}
\begin{table}
\centering
\begin{tabular}{lcc}
\toprule
Ítem & Conteo & Orden \\
\midrule
pan & 4 & 1 \\
leche & 4 & 2 \\
huevos & 3 & 3 \\
mantequilla & 3 & 4 \\
cereal & 1 & eliminado \\
\bottomrule
\end{tabular}
\end{table}

\column{0.54\textwidth}
\begin{table}
\centering
\begin{tabular}{lll}
\toprule
ID & Original & Ordenada sin cereal \\
\midrule
T1 & pan, leche, huevos & pan, leche, huevos \\
T2 & pan, mantequilla & pan, mantequilla \\
T3 & leche, huevos, cereal & leche, huevos \\
T4 & pan, leche, mantequilla & pan, leche, mantequilla \\
T5 & pan, leche, huevos, mantequilla & pan, leche, huevos, mantequilla \\
\bottomrule
\end{tabular}
\end{table}
\end{columns}
\end{frame}

%------------------------------------------------
\begin{frame}[fragile]{FP-Growth a mano II: construir el FP-tree}
\scriptsize
\begin{columns}[T]
\column{0.45\textwidth}
\textbf{Inserción resumida}
\begin{enumerate}
  \item T1 crea la ruta pan-leche-huevos.
  \item T2 comparte pan y agrega mantequilla bajo pan.
  \item T3 crea una ruta nueva leche-huevos desde la raíz.
  \item T4 comparte pan-leche y agrega mantequilla.
  \item T5 comparte pan-leche-huevos y agrega mantequilla.
\end{enumerate}

\column{0.50\textwidth}
\textbf{FP-tree resultante}
\begin{verbatim}
raiz
|- pan:4
|  |- leche:3
|  |  |- huevos:2
|  |  |  `- mantequilla:1
|  |  `- mantequilla:1
|  `- mantequilla:1
`- leche:1
   `- huevos:1
\end{verbatim}
\end{columns}
\end{frame}

%------------------------------------------------
\begin{frame}{FP-Growth a mano III: bases condicionales y patrones}
\scriptsize
\begin{columns}[T]
\column{0.47\textwidth}
\textbf{Sufijo: mantequilla}
\begin{table}
\centering
\begin{tabular}{lc}
\toprule
Prefijo condicional & Conteo \\
\midrule
pan & 1 \\
pan, leche & 1 \\
pan, leche, huevos & 1 \\
\bottomrule
\end{tabular}
\end{table}
Conteos condicionales:
\[
pan=3,\quad leche=2,\quad huevos=1
\]
Patrones frecuentes:
\[
\{pan,mantequilla\}:3,\quad \{leche,mantequilla\}:2
\]
\[
\{pan,leche,mantequilla\}:2
\]

\column{0.48\textwidth}
\textbf{Sufijo: huevos}
\begin{table}
\centering
\begin{tabular}{lc}
\toprule
Prefijo condicional & Conteo \\
\midrule
pan, leche & 2 \\
leche & 1 \\
\bottomrule
\end{tabular}
\end{table}
Conteos condicionales:
\[
leche=3,\quad pan=2
\]
Patrones frecuentes:
\[
\{leche,huevos\}:3,\quad \{pan,huevos\}:2
\]
\[
\{pan,leche,huevos\}:2
\]
\end{columns}
\end{frame}

%------------------------------------------------
\begin{frame}{Resultado manual común de los tres algoritmos}
\scriptsize
\begin{table}
\centering
\begin{tabular}{lcc}
\toprule
Itemset frecuente & Conteo & Soporte \\
\midrule
\{pan\}, \{leche\} & 4 & 0.80 \\
\{huevos\}, \{mantequilla\} & 3 & 0.60 \\
\{pan, leche\} & 3 & 0.60 \\
\{pan, huevos\} & 2 & 0.40 \\
\{pan, mantequilla\} & 3 & 0.60 \\
\{leche, huevos\} & 3 & 0.60 \\
\{leche, mantequilla\} & 2 & 0.40 \\
\{pan, leche, huevos\} & 2 & 0.40 \\
\{pan, leche, mantequilla\} & 2 & 0.40 \\
\bottomrule
\end{tabular}
\end{table}

\begin{block}{Conclusión}
Apriori, Eclat y FP-Growth recorren el espacio de búsqueda de forma distinta, pero con el mismo soporte mínimo recuperan los mismos itemsets frecuentes.
\end{block}
\end{frame}



%------------------------------------------------
\begin{frame}[fragile]{R: paquetes necesarios}
\begin{lstlisting}[style=Rstyle]
# Instalar una vez
install.packages("arules")
install.packages("arulesViz")
install.packages("tidyverse")

# Cargar en cada sesion
library(arules)
library(arulesViz)
library(tidyverse)
\end{lstlisting}

\begin{block}{Paquete central}
\texttt{arules} implementa estructuras de transacciones, Apriori, Eclat y funciones para inspeccionar reglas.
\end{block}
\end{frame}

%------------------------------------------------
\begin{frame}[fragile]{R: cargar datos Groceries}
\begin{lstlisting}[style=Rstyle]
library(arules)
data("Groceries")

Groceries
summary(Groceries)
itemFrequency(Groceries)[1:10]
\end{lstlisting}

\begin{itemize}
  \item \texttt{Groceries} es un dataset transaccional clásico incluido en \texttt{arules}.
  \item Cada transacción corresponde a una compra.
  \item Cada ítem corresponde a un producto comprado.
\end{itemize}
\end{frame}

%------------------------------------------------
\begin{frame}[fragile]{R: explorar frecuencia de ítems}
\begin{lstlisting}[style=Rstyle]
# Frecuencia relativa de items
freq <- itemFrequency(Groceries, type = "relative")
head(sort(freq, decreasing = TRUE), 10)

# Grafico de los 20 items mas frecuentes
itemFrequencyPlot(
  Groceries,
  topN = 20,
  type = "relative",
  main = "Items mas frecuentes"
)
\end{lstlisting}

\begin{block}{Interpretación}
Antes de generar reglas conviene revisar cuáles ítems dominan el dataset, porque pueden producir reglas triviales.
\end{block}
\end{frame}

%------------------------------------------------
\begin{frame}[fragile]{R: generar reglas con Apriori}
\begin{lstlisting}[style=Rstyle]
rules <- apriori(
  Groceries,
  parameter = list(
    supp = 0.01,
    conf = 0.30,
    minlen = 2
  )
)

rules
summary(rules)
inspect(head(rules, 10))
\end{lstlisting}

\begin{itemize}
  \item \texttt{supp = 0.01}: la regla debe aparecer en al menos 1\% de las compras.
  \item \texttt{conf = 0.30}: al menos 30\% de confianza.
  \item \texttt{minlen = 2}: antecedente y consecuente no vacíos.
\end{itemize}
\end{frame}

%%------------------------------------------------
%\begin{frame}[fragile]{R: ordenar por lift}
%\begin{lstlisting}[style=Rstyle]
%rules_lift <- sort(rules, by = "lift", decreasing = TRUE)
%inspect(head(rules_lift, 10))
%
%# Extraer medidas de calidad
%quality(head(rules_lift, 5))
%\end{lstlisting}
%
%\begin{block}{Buena práctica}
%No seleccionar reglas solo por confianza. Revisar también soporte, lift, tamaño de la regla y plausibilidad del patrón.
%\end{block}
%\end{frame}
%
%%------------------------------------------------
%\begin{frame}[fragile]{R: filtrar reglas con un consecuente específico}
%\begin{lstlisting}[style=Rstyle]
%# Reglas que predicen compra de whole milk
%rules_milk <- subset(
%  rules,
%  subset = rhs %in% "whole milk" & lift > 1.2
%)
%
%rules_milk <- sort(rules_milk, by = "confidence")
%inspect(head(rules_milk, 10))
%\end{lstlisting}
%
%\begin{exampleblock}{Pregunta analítica}
%¿Qué productos son buenos predictores de compra de leche? ¿La regla es útil comercialmente o solo refleja productos muy frecuentes?
%\end{exampleblock}
%\end{frame}
%
%%------------------------------------------------
%\begin{frame}[fragile]{R: eliminar reglas redundantes}
%\begin{lstlisting}[style=Rstyle]
%# Reglas redundantes: una regla mas especifica
%# no mejora la calidad de una regla mas general
%redundant <- is.redundant(rules_lift)
%sum(redundant)
%
%rules_pruned <- rules_lift[!redundant]
%inspect(head(rules_pruned, 10))
%\end{lstlisting}
%
%\begin{block}{Interpretación}
%La poda facilita leer el resultado y evita reportar variantes casi idénticas de la misma asociación.
%\end{block}
%\end{frame}
%
%%------------------------------------------------
%\begin{frame}[fragile]{R: visualización de reglas}
%\begin{lstlisting}[style=Rstyle]
%library(arulesViz)
%
%# Grafico soporte vs confianza, coloreado por lift
%plot(rules_pruned, method = "scatterplot")
%
%# Red de reglas principales
%plot(head(rules_pruned, 20), method = "graph")
%
%# Matriz agrupada
%plot(head(rules_pruned, 50), method = "grouped")
%\end{lstlisting}
%
%\begin{itemize}
%  \item Los gráficos ayudan a detectar reglas fuertes, redundantes o aisladas.
%  \item En informes, conviene mostrar pocas reglas bien seleccionadas.
%\end{itemize}
%\end{frame}
%
%%------------------------------------------------
%\begin{frame}[fragile]{R: usar Eclat para itemsets frecuentes}
%\begin{lstlisting}[style=Rstyle]
%itemsets <- eclat(
%  Groceries,
%  parameter = list(supp = 0.02, maxlen = 4)
%)
%
%itemsets <- sort(itemsets, by = "support")
%inspect(head(itemsets, 10))
%\end{lstlisting}
%
%\begin{block}{Diferencia práctica}
%\texttt{eclat()} retorna itemsets frecuentes. Para obtener reglas, normalmente se generan reglas desde itemsets frecuentes o se usa directamente \texttt{apriori()}.
%\end{block}
%\end{frame}
%
%%------------------------------------------------
%\begin{frame}[fragile]{R: crear transacciones desde una tabla}
%\begin{lstlisting}[style=Rstyle]
%compras <- tibble::tribble(
%  ~id, ~producto,
%  1, "pan", 1, "leche", 1, "huevos",
%  2, "pan", 2, "mantequilla",
%  3, "leche", 3, "huevos", 3, "cereal",
%  4, "pan", 4, "leche", 4, "mantequilla",
%  5, "pan", 5, "leche", 5, "huevos", 5, "mantequilla"
%)
%
%trans <- as(split(compras$producto, compras$id), "transactions")
%inspect(trans)
%\end{lstlisting}
%
%\begin{block}{Idea}
%Muchos datos reales están en formato largo: una fila por transacción e ítem. Hay que convertirlos a \texttt{transactions}.
%\end{block}
%\end{frame}
%
%%------------------------------------------------
%\begin{frame}[fragile]{R: reglas en el ejemplo pequeño}
%\begin{lstlisting}[style=Rstyle]
%rules_small <- apriori(
%  trans,
%  parameter = list(supp = 0.4, conf = 0.6, minlen = 2)
%)
%
%inspect(sort(rules_small, by = "lift"))
%\end{lstlisting}
%
%\begin{alertblock}{Advertencia}
%Con muestras pequeñas, las reglas pueden parecer fuertes por azar. Es importante reportar el tamaño muestral y validar en nuevos datos.
%\end{alertblock}
%\end{frame}
%
%%------------------------------------------------
%\begin{frame}[fragile]{Datos no transaccionales: discretización}
%\begin{itemize}
%  \item Las reglas de asociación requieren variables discretas o binarias.
%  \item Variables continuas deben discretizarse: bajo/medio/alto, cuartiles, umbrales clínicos, etc.
%  \item La discretización afecta las reglas; debe justificarse.
%\end{itemize}
%
%\begin{lstlisting}[style=Rstyle]
%# Ejemplo con iris: convertir variables continuas a categorias
%iris_disc <- discretizeDF(
%  iris,
%  methods = list(
%    Sepal.Length = "frequency",
%    Sepal.Width  = "frequency",
%    Petal.Length = "frequency",
%    Petal.Width  = "frequency"
%  )
%)
%head(iris_disc)
%\end{lstlisting}
%\end{frame}
%
%%------------------------------------------------
%\begin{frame}[fragile]{R: reglas desde iris discretizado}
%\begin{lstlisting}[style=Rstyle]
%trans_iris <- as(iris_disc, "transactions")
%
%rules_iris <- apriori(
%  trans_iris,
%  parameter = list(supp = 0.10, conf = 0.80, minlen = 2),
%  appearance = list(rhs = c("Species=setosa",
%                            "Species=versicolor",
%                            "Species=virginica"),
%                    default = "lhs")
%)
%
%inspect(sort(rules_iris, by = "lift"))
%\end{lstlisting}
%
%\begin{block}{Uso}
%Esta estrategia permite encontrar combinaciones de rangos de variables asociadas a una clase.
%\end{block}
%\end{frame}
%
%%------------------------------------------------
%\begin{frame}{Interpretación de una regla}
%Supongamos una regla:
%\[
%\{Petal.Length = alto, Petal.Width = alto\} \Rightarrow \{Species = virginica\}
%\]
%
%\begin{itemize}
%  \item \textbf{Soporte}: ¿en cuántas muestras aparece este patrón completo?
%  \item \textbf{Confianza}: ¿qué proporción de flores con pétalos altos son virginica?
%  \item \textbf{Lift}: ¿cuánto más probable es virginica bajo el antecedente frente a su frecuencia basal?
%  \item \textbf{Utilidad}: ¿la regla aporta una explicación simple o solo repite una relación obvia?
%\end{itemize}
%\end{frame}
%
%%------------------------------------------------
%\begin{frame}[fragile]{Asociación y prueba estadística}
%\begin{lstlisting}[style=Rstyle]
%# Tomar una regla y construir tabla 2x2
%r <- rules_pruned[1]
%
%# Matriz de incidencia de transacciones
%X <- is.subset(lhs(r), Groceries)
%Y <- is.subset(rhs(r), Groceries)
%
%tab <- table(antecedente = as.vector(X), consecuente = as.vector(Y))
%tab
%
%fisher.test(tab)
%chisq.test(tab)
%\end{lstlisting}
%
%\begin{alertblock}{Cuidado}
%Si se prueban miles de reglas, hay problema de comparaciones múltiples. No basta con un p-valor aislado.
%\end{alertblock}
%\end{frame}

%------------------------------------------------
\begin{frame}{Selección responsable de reglas}
\begin{enumerate}
  \item Definir umbrales razonables de soporte y confianza.
  \item Priorizar reglas con lift alto, pero no extremadamente raras.
  \item Eliminar reglas redundantes o triviales.
  \item Validar reglas en datos independientes, cuando sea posible.
  \item Incorporar conocimiento del dominio antes de concluir causalidad o acción.
\end{enumerate}

\begin{block}{Recordar}
Las reglas de asociación descubren patrones, no prueban causalidad.
\end{block}
\end{frame}

%------------------------------------------------
\begin{frame}{Limitaciones frecuentes}
\begin{itemize}
  \item Explosión combinatoria cuando hay muchos ítems y soporte bajo.
  \item Reglas triviales por ítems muy frecuentes.
  \item Sensibilidad a la discretización de variables continuas.
  \item Dificultad para interpretar miles de reglas simultáneamente.
  \item No modela ajuste por confusores de forma directa.
  \item Riesgo de sobreinterpretar correlaciones como causalidad.
\end{itemize}
\end{frame}

%------------------------------------------------
\begin{frame}{Extensiones y variantes}
\begin{itemize}
  \item Reglas secuenciales: consideran orden temporal de eventos.
  \item Reglas ponderadas: algunos ítems tienen mayor importancia o valor.
  \item Reglas cuantitativas: incorporan variables numéricas discretizadas.
  \item Reglas negativas: identifican ausencia conjunta o exclusión.
  \item Reglas de clasificación: consecuente restringido a una clase objetivo.
\end{itemize}

\begin{exampleblock}{Ejemplo aplicado}
En salud, se puede restringir el consecuente a ``respuesta a tratamiento'' y buscar combinaciones de biomarcadores asociadas a respuesta o no respuesta.
\end{exampleblock}
\end{frame}

%------------------------------------------------
\begin{frame}{Ejercicio 1: canasta de compras}
\begin{enumerate}
  \item Cargue el dataset \texttt{Groceries}.
  \item Explore los 20 ítems más frecuentes.
  \item Genere reglas con \texttt{supp = 0.005} y \texttt{conf = 0.25}.
  \item Ordene por lift y seleccione 10 reglas interpretables.
  \item Elimine reglas redundantes.
  \item Discuta cuál regla sería útil para una estrategia comercial.
\end{enumerate}
\end{frame}

%%------------------------------------------------
%\begin{frame}{Ejercicio 2: reglas asociadas a clase}
%\begin{enumerate}
%  \item Use el dataset \texttt{iris}.
%  \item Discretice las variables numéricas en tres categorías.
%  \item Transforme el dataset a formato \texttt{transactions}.
%  \item Genere reglas cuyo consecuente sea la especie.
%  \item Compare reglas para \texttt{setosa}, \texttt{versicolor} y \texttt{virginica}.
%  \item Interprete soporte, confianza y lift de las mejores reglas.
%\end{enumerate}
%\end{frame}

%%------------------------------------------------
%\begin{frame}{Ejercicio 3: interpretación crítica}
%Para una regla con alto lift pero bajo soporte:
%
%\begin{enumerate}
%  \item ¿Por qué podría ser interesante?
%  \item ¿Por qué podría ser inestable?
%  \item ¿Cómo la validaría en nuevos datos?
%  \item ¿Qué criterio usaría para decidir si reportarla?
%\end{enumerate}
%
%\begin{block}{Discusión}
%En minería de datos, no siempre buscamos el patrón más fuerte, sino el patrón más útil, estable e interpretable.
%\end{block}
%\end{frame}

%------------------------------------------------
\begin{frame}{Resumen}
\begin{itemize}
  \item El análisis de asociación busca relaciones frecuentes del tipo $X \Rightarrow Y$.
  \item Soporte, confianza y lift responden preguntas distintas.
  \item Apriori usa poda antimonótona; Eclat usa formato vertical; FP-Growth usa un árbol comprimido.
  \item En R, \texttt{arules} permite generar, filtrar, ordenar y visualizar reglas.
  \item La interpretación debe considerar redundancia, frecuencia basal, estabilidad y conocimiento del dominio.
\end{itemize}
\end{frame}

%------------------------------------------------
\begin{frame}{Referencias}
\begin{itemize}
  \item Agrawal, R. y Srikant, R. (1994). Fast algorithms for mining association rules.
  \item Han, J., Pei, J. y Yin, Y. (2000). Mining frequent patterns without candidate generation.
  \item Hahsler, M., Grün, B. y Hornik, K. (2005). \texttt{arules}: Mining association rules and frequent itemsets.
  \item Hahsler, M. (2017). \texttt{arulesViz}: Interactive visualization of association rules.
\end{itemize}
\end{frame}

\end{document}
